% IEEE TIFS submission format: two-column journal, single-anonymized review.
\documentclass[10pt,journal]{IEEEtran}
\providecommand{\Description}[1]{}
\setlength{\headheight}{20.5pt}
% Keep this manuscript standalone: if the appendix is mentioned, use plain text
% such as "Appendix" instead of cross-document \ref commands to appendix.tex.

\usepackage{qtree}
\usepackage{graphicx}
\usepackage{listings}
\usepackage{xcolor}
\usepackage{amsfonts}
\usepackage{mathrsfs}
\usepackage{amsmath}
\usepackage{amsthm}
\usepackage{amscd}
\usepackage{tabularx}
\usepackage{multirow}
\usepackage{float}
% No subfigure package is needed in the compact IEEE main paper.
\usepackage{array}
\usepackage{algorithm}
\usepackage{bm}
\usepackage{comment}
\usepackage{makecell}
% IEEEtran supplies the conference caption formatting.
\newcommand*{\defeq}{\stackrel{\text{def}}{=}}

\usepackage{algpseudocode}
\usepackage[hidelinks]{hyperref}
\raggedbottom

\algnewcommand\algorithmicforeach{\textbf{for each}}
\algdef{S}[FOR]{ForEach}[1]{\algorithmicforeach\ #1\ \algorithmicdo}

\newtheorem{definition}{Definition}
\newtheorem{theorem}{Theorem}
\newtheorem{lemma}{Lemma}
\newtheorem{proposition}{Proposition}
\newtheorem{corollary}{Corollary}
\newtheorem{assumption}{Assumption}

\title{Not Where, but What: Operator Choice in Spatially Selective Visual Location Privacy}
% TIFS uses single-anonymized review: reviewers are anonymous, authors are not,
% so the author block is disclosed rather than withheld.
\author{Bo~Ma,
        Weiqi~Yan,
        and~Jinsong~Wu,~\IEEEmembership{Senior~Member,~IEEE}
\thanks{B.~Ma is with Resideo Technologies Inc., Scottsdale, AZ, USA, and with
the Auckland University of Technology, Auckland, New Zealand (e-mail:
rcn4743@aut.ac.nz).}
\thanks{W.~Yan is with the Auckland University of Technology, Auckland, New
Zealand.}
\thanks{J.~Wu is with the Guilin University of Electronic Technology, Guilin,
China.}}
\begin{document}
\maketitle
\begin{abstract}
Visual location-privacy mechanisms are built around a single design decision: predict which pixels carry the location signal, then concentrate a limited perturbation budget there. The budget's \emph{placement} is what these systems learn, tune and report; what the budget is spent \emph{on} is almost always fixed at additive isotropic Gaussian noise and never treated as a variable. We show this is the wrong way round. We give an evaluation protocol that controls both axes --- redistributing a mechanism's own weights uniformly, cyclically or by pixel permutation under a verified energy-conservation gate, and comparing perturbation operators at matched \emph{delivered} distortion rather than nominal budget --- and apply it to a DCRF-guided mechanism built as a testbed, on a real place-labelled MSLS benchmark. Placement is inert. Across eight placement rules, three rules driven by published segmentation models, an attacker-gradient oracle, the strictly correct margin-gradient oracle, two checkpoints, six attacker backbones and four operators, nothing beats spreading the budget uniformly at the quality-preserving operating point; the protocol also exposes that the testbed's own learned map is numerically constant, so it performs no spatial selection at all while still producing a plausible privacy--utility curve. The operator is not inert. Once the budget is large enough to move retrieval at all, changing only what it is spent on moves Top-1 from $0.159$ to $0.053$ at identical delivered distortion and identical uniform placement, and the placement effect that finally appears reverses sign between operators, so a rule tuned on one is harmful on another. A controlled retrieval task with an exactly known Jacobian explains why: at matched energy, three operators differing in spatial structure are interchangeable, while one differing in \emph{directional} control drives Top-1 to $0.001$ under uniform placement and no selectivity whatsoever. Spatial allocation --- of the map or of the operator --- is not where visual location privacy comes from. We recommend energy-matched placement controls and matched-distortion operator comparison as standard practice.
% \vspace{1em}
% \noindent\textbf{Code:} \url{https://github.com/mabo1215/PPEDCRF}
\end{abstract}

\section{Introduction}
Mobile-recorded videos --- from dashcams, ADAS platforms, bodycams, drones and
wearable cameras --- are routinely uploaded for safety investigation, incident
auditing and model improvement. Stripping explicit metadata such as GPS
coordinates is necessary but not sufficient: the background of a single frame
may contain distinctive landmarks matchable against large-scale geo-tagged
databases. This paper addresses \emph{location privacy} under such a
background-based retrieval attacker, a setting distinct from face or
license-plate anonymization because buildings, road layouts, signage and
skylines expose where a video was recorded even when every identity is
hidden.

The dominant response is \emph{selective} perturbation: predict which pixels carry the sensitive signal, then concentrate a limited noise budget there, so that quality is spent only where privacy requires it. Such methods are normally validated by showing a better privacy--utility curve than perturbing the whole frame. We observe that this comparison cannot separate two very different explanations of the same curve. A selective mechanism might work because it identified the right pixels, or merely because a spatially concentrated perturbation of a given total energy happens to disturb an embedding as much as a diffuse one --- in which case the learned map is decorative, and any map of equal energy would do.

Distinguishing these requires a control that holds perturbation energy fixed and varies only \emph{where} it lands. To our knowledge such controls are not standard practice in this literature.

There is a second variable, and the field holds it fixed without comment. A
selective mechanism decides where the budget goes, but also what it is spent
\emph{on}, and that choice is almost universally additive isotropic Gaussian
noise: placement is learned, tuned and ablated, while the operator is
inherited. If the leverage lies in the operator rather than the map, no amount
of evidence about placement would reveal it, because placement studies never
vary the operator.

We therefore ask a methodological question rather than a design one: \textbf{which of the two axes a selective mechanism controls --- where the budget lands, or what it is spent on --- actually determines the privacy it delivers?}

\paragraph{Threat model and released object.}
Let $q$ denote a released video frame and $\mathcal{D}=\{(s_i,\ell_i)\}_{i=1}^N$ a gallery of geo-tagged reference images with location labels $\ell_i$. The attacker computes an embedding $f(\cdot)$ and ranks gallery images by similarity $\mathrm{sim}(f(q),f(s_i))$. A retrieval attack succeeds if the correct location appears in the Top-$k$ list. Our defense modifies the released frame before the attacker sees it. The downstream utility target is a detector or segmenter that still operates on the sanitized frame. A short neighboring-frame context may be used internally to stabilize the support estimate, but the primary released object and evaluated attack are frame-level.

Encryption does not address this setting: once a decrypted frame reaches a
perception module for auditing, training or analytics, its background cues are
equally available to a retrieval attacker. The objective is
utility-preserving sanitization.

\paragraph{Scope of the claims.}
This paper makes an evaluation claim, not a superiority claim. The mechanism we test uses a Gaussian-inspired calibration index to select an image-space noise scale; we do \emph{not} claim a formal $(\varepsilon,\delta)$-DP guarantee, and the index is an engineering control for choosing an operating point rather than a certified budget. We also do not claim that spatial selectivity is useless in general --- our own budget sweep shows it is not. What we establish is narrower and checkable. First, over eight placement rules and two checkpoints in the core study (98 paired comparisons over six attacker backbones), plus 23 further paired tests on the real place-labelled benchmark covering three published-model placements, the margin-gradient oracle and four perturbation operators, nothing outperforms uniform spreading \emph{at the quality-preserving operating point}, and the one map with real spatial selectivity performs significantly worse. Second, the operator is a live design variable where placement is not, but
only once the budget is large enough to move retrieval at all; at the
operating point every operator we tested leaves accuracy essentially where
uniform Gaussian noise leaves it. We do \emph{not} claim a well-chosen
operator solves the problem at quality-preserving budgets --- none of ours
does. The claim is comparative. A different mechanism or task could behave otherwise; the protocol is offered so that such a claim can be checked rather than assumed.

\begin{comment}
\begin{figure*}[!t]
  \centering
  \includegraphics[width=\textwidth]{figs/twopic_clear.png}
  \caption{A released video frame can be matched to a geo-tagged reference image through background cues even when explicit metadata are removed.}
  \Description{A two-panel example showing a released video frame and a visually similar geo-tagged reference image, illustrating that background structures alone can reveal location.}
  \label{fig:twopic}
\end{figure*}
\end{comment}

\subsection{Related Work}
\textbf{Visual anonymization for multimedia.}
Visual privacy has been surveyed extensively~\cite{padilla2015visual,ribaric2016deidentification}; classical methods obfuscate foreground entities---mainly faces and license plates---by blurring, pixelation, or cryptographic protocols~\cite{Chamikara2020101951,Li2019212,avidan2007efficient,erkin2009privacy}. Generative approaches instead synthesize de-identified faces: DeepPrivacy~\cite{hukkelaas2019deepprivacy} replaces faces via a pose-conditioned GAN, CIAGAN~\cite{maximov2020ciagan} performs identity swaps preserving expression and pose, and head inpainting~\cite{sun2018natural} avoids the artifacts of na\"{\i}ve blurring. Oh et al.~\cite{oh2016faceless} showed recognition networks can re-identify people even with faces fully occluded, and Zhou and Pun~\cite{zhou2020personal} pixelate bystanders in live streams. All are foreground-centric: none address a sensitive signal residing in the \emph{background scene} that supports geo-localization, the threat model targeted here.

\textbf{Privacy-preserving perception pipelines.}
A parallel line modifies inputs, training, or intermediate representations to limit leakage: privacy-preserving medical detection~\cite{liu2019privacy}, learned face anonymization that preserves action detection~\cite{ren2018learning}, federated and encrypted inference~\cite{kaissis2020secure}, and GANobfuscator~\cite{xu2019ganobfuscator}, which perturbs images before sharing to block attribute inference. These emphasize downstream utility but again assume a foreground-centric risk, whereas our primary risk is the scene context enabling geo-localization.

\textbf{Visual place recognition and geo-localization.}
VPR matches a query to geo-tagged references using learned
descriptors~\cite{lowry2016visual,masone2021survey}, from trainable VLAD
pooling over deep backbones~\cite{arandjelovic2016netvlad,he2016deep} and
day--night robustness~\cite{torii201524} to multi-scale patch
fusion~\cite{hausler2021patchnetvlad} and the stronger
CosPlace~\cite{berton2022rethinking} and MixVPR~\cite{ali2023mixvpr}. These
define the attacker side of our threat model: the stronger the embedding, the
harder any attacker-agnostic defense becomes. We evaluate three dedicated VPR
embedders alongside classification-style attackers, but not the full design
space, so ours is bounded transfer evidence rather than universal robustness.

\textbf{Scene-level and location privacy.}
Point-cloud localization representations can leak detailed scene
content~\cite{pittaluga2019revealing}, motivating localization that hides
geometric structure from the server~\cite{speciale2019privacy};
geo-indistinguishability~\cite{andres2013geo} extends DP to location-based
services and Shokri et al.~\cite{shokri2011quantifying} quantify location
privacy against inference attacks. These protect geometric or coordinate-level
representations, whereas the setting here is the released pixel-domain frame
against retrieval-based geo-localization.

\textbf{Differential privacy and perturbation-based defenses.}
DP bounds leakage formally~\cite{dwork2014algorithmic} through the calibrated
Gaussian mechanism~\cite{balle2018improving}, R\'{e}nyi
composition~\cite{mironov2017renyi} and DP-SGD~\cite{abadi2016deep}; in
multimedia, VideoDP~\cite{wang2020videodp} applies DP to video analytics and
Roy and Boddeti~\cite{roy2019mitigating} limit representation leakage by
maximum entropy. Perturbation defenses are also studied
adversarially~\cite{goodfellow2015explaining,madry2018towards}. Our testbed
uses a Gaussian-inspired calibration index as a strength knob rather than
claiming end-to-end DP.

\textbf{Conditional random fields for spatial--temporal inference.}
Efficient mean-field inference for fully connected
CRFs~\cite{krahenbuhl2011efficient}, CRF-as-RNN~\cite{zheng2015conditional}
and dense CRFs inside segmentation~\cite{chen2017deeplab} established the
spatial machinery, and dynamic CRFs~\cite{wang2005dynamic} add temporal edges
across frames. Our testbed combines mean-field-style spatial smoothing with
temporal consistency to produce the support maps that guide its perturbation.

\textbf{Privacy--utility trade-off optimization.}
A growing body of work studies the privacy--utility tension in image and video processing~\cite{wu2018towards}. McPherson et al.~\cite{mcpherson2016defeating} showed that mosaicing and blurring can be inverted by trained networks, motivating more principled perturbation. Across this literature, spatial selectivity is typically validated by a favourable trade-off curve against a non-selective baseline. Our contribution is orthogonal to any particular design: we supply the missing control that holds perturbation energy fixed while varying only its placement, and report what it shows for a mechanism of this family.


\subsection{Contributions}
Our contributions are:
\begin{itemize}
    \item \textbf{A protocol that controls both design axes.} Given any mechanism emitting per-pixel perturbation weights, we construct three placement controls --- uniform, cyclically shifted and pixel-permuted --- preserving the weights' sum of squares exactly under an energy-conservation gate, and compare retrieval pairwise on identical queries and seeds. We then add the axis this literature leaves fixed, comparing perturbation \emph{operators} at matched \emph{delivered} MSE, solved per frame through the pixel clamp. Together these separate ``the placement matters'' from ``the energy matters'' from ``the operator matters,'' and matching delivered rather than nominal distortion removes a clipping confound that nominal-energy matching can only bound.
	\item \textbf{Negative results obtained with that protocol.} Applied to PPEDCRF --- a DCRF-guided selective mechanism we implement as a testbed --- the protocol immediately exposes that its released map is a constant, so the mechanism performs no spatial selection at all while still producing a plausible privacy--utility curve. Across eight energy-matched placements on two checkpoints (one constant, one genuinely selective), including an attacker-gradient oracle, no placement beats uniform spreading at the operating point in 98 comparisons, and the genuinely selective map is significantly worse.
	\item \textbf{The operator carries the leverage that placement does not.} Holding placement fixed and varying what the budget is spent on --- additive Gaussian, spatially correlated, selective low-pass, selective block quantisation --- at matched delivered MSE, the null survives every operator at the operating point. But once the budget is large enough to move retrieval at all, the operator alone moves Top-1 from 0.159 to 0.053 at identical distortion and identical uniform placement, and the sign of the placement effect reverses between operators, so a rule tuned on one operator is harmful on another. In the controlled model, three operators differing in spatial structure are interchangeable at matched energy whereas one differing in \emph{directional} control drives Top-1 to 0.001 under uniform placement and no selectivity at all. Spatial allocation, whether by the map or by the operator, is not the source of the protection.
	\item \textbf{A mechanistic account, checked against a known Jacobian.} Attacker sensitivity is in fact strongly concentrated (top decile carries 67.7\% of gradient energy), so the null is not explained by homogeneity; rather, gradient magnitude fixes \emph{where} noise goes but not \emph{which direction} it moves the embedding, which is why an optimized white-box perturbation at the same quality succeeds where every placement fails. We test this in a controlled retrieval task whose encoder is known in closed form, so sensitivity is exact rather than estimated. There the first-order prediction is confirmed to within 4\% across 168 cells and the oracle advantage it predicts is large, establishing that the measured null is not an artifact of the algebra; switching on a bounded output range and realistic task difficulty then removes that advantage, while a nonlinear encoder does not, isolating the two causes. The account predicts, and a budget sweep confirms, that placement becomes exploitable only once the perturbation leaves the small-amplitude regime --- that is, only after the image quality these mechanisms exist to preserve has already been spent. The results hold across six attacker backbones, three gallery sizes, a real GPS-labeled MSLS benchmark spanning eight cities, and both threat models, with cluster-aware significance throughout and a one-command artifact that recomputes every reported number from raw outputs.
\end{itemize}

\section{The Mechanism Under Test}
\label{sec:method}

To exercise the protocol we need a concrete spatially selective mechanism that
is representative of the family described above: one that predicts where the
sensitive signal lies, concentrates a fixed budget there, and exposes the
per-pixel weights the protocol needs. We build PPEDCRF for this purpose. It is
deliberately conventional --- a learned support map, temporal smoothing, a
normalized amplitude rule, and calibrated Gaussian noise --- because the
question is not whether this particular design is the best of its kind, but
whether a design of this kind derives its protection from placement at all. A
mechanism of our own construction also lets us report a negative result
without misrepresenting anyone else's system.

\subsection{Pipeline Overview}
The mechanism operates on a frame stream $\mathcal{X}=\{I_t\}_{t=1}^T$,
following $I_t \rightarrow u_t \rightarrow p_t \rightarrow \alpha_t
\rightarrow I_t'$: a unary support logit $u_t$, a DCRF-refined map $p_t$ that
may reuse the previous frame as a temporal prior, an NCP control map
$\alpha_t$ setting perturbation strength, and the released frame $I_t'$ the
attacker observes. Only the product $w_t = \alpha_t \odot p_t$ matters for
what follows: it is the per-pixel weight vector our protocol redistributes.

\begin{comment}
\begin{figure*}[!t]
    \centering
    \includegraphics[width=0.8\textwidth]{figs/architecture_of_solution.png}
    \caption{Overview of the PPEDCRF pipeline. A unary candidate-support predictor produces per-frame logits, DCRF enforces spatial--temporal consistency, NCP rescales perturbation strength, and indexed Gaussian noise is injected only in the selected background support before release.}
    \Description{A pipeline diagram with four stages: unary background-sensitivity prediction, dynamic CRF refinement across neighboring frames, normalized control penalty scaling, and selective Gaussian perturbation of the released frame.}
    \label{fig:arc}
\end{figure*}
\end{comment}

\begin{comment}
\begin{table}[t]
\centering
\caption{Compact notation summary.}
\label{tab:symbols}
\begin{tabularx}{\linewidth}{>{\raggedright\arraybackslash}p{0.16\linewidth} X}
\hline
Symbol & Meaning \\
\hline
$I_t$ & Input frame at time $t$. \\
$u_t$ & Unary candidate-support logit predicted from $I_t$. \\
$p_t$ & DCRF-refined continuous support map. \\
$\alpha_t$ & NCP control map that rescales perturbation strength. \\
$\sigma_0$ & Base indexed-Gaussian noise scale selected by the calibration index. \\
$I_t'$ & Sanitized frame released to the attacker or downstream model. \\
$f(\cdot)$ & External retrieval backbone used by the attacker. \\
$\mathcal{D}$ & Retrieval gallery of candidate reference images. \\
\hline
\end{tabularx}
\end{table}
\end{comment}

\subsection{DCRF Mask Inference}
The unary network $g_\theta(\cdot)$ predicts a per-pixel candidate background-support logit
\begin{equation}
\label{eq:unary}
u_t = g_\theta(I_t).
\end{equation}
Here $\sigma(z)=1/(1+\exp(-z))$ is the element-wise logistic sigmoid. The
released unary predictor is a newly implemented lightweight encoder--decoder,
not an adapted VPR or segmentation backbone but a small encoder--decoder of
87,441 parameters, trained with binary cross-entropy against pixel-level
support masks when these are available. The monitoring pool used for retrieval
has no independent pixel-level support ground truth, so we do not interpret
the unary output as a validated privacy-sensitivity estimate.

We initialize the continuous candidate-support map as $p_t^{(0)}=\sigma(u_t)$ and refine it with a dynamic CRF using a small number of mean-field iterations that matches the released implementation:
\begin{equation}
\label{eq:dcrf}
p_t^{(k+1)} =
\sigma\!\left(
u_t + \lambda_s \big(S(p_t^{(k)}) - p_t^{(k)}\big)
      + \lambda_\tau \big(\hat{p}_{t-1} - p_t^{(k)}\big)
\right),
\end{equation}
where $S(\cdot)$ is average-pooling-based spatial smoothing with a
$3\times3$ kernel, stride 1, and reflect padding in the released
implementation; $\hat{p}_{t-1}$ is the previous refined map warped into frame
$t$ (or reused directly when optical flow is unavailable), and $\lambda_s,
\lambda_\tau$ are spatial and temporal weights. We use five mean-field-style
iterations, $\lambda_s=\lambda_\tau=2$, and no optical-flow warp in the
reported controlled benchmark. After these fixed iterations, the final
refined mask is denoted by $p_t$.


\subsection{NCP Control and Gaussian-Inspired Calibration}
The normalized control penalty (NCP) rescales the refined mask rather than redefining the threat model. In the released code, NCP is implemented as a normalized control map
\begin{equation}
\label{eq:ncp}
\alpha_t = \alpha \cdot \frac{p_t}{\max(p_t)+\epsilon},
\end{equation}
where $\alpha$ is a global control knob and $\epsilon$ avoids division by zero. Optional class-specific weights can be multiplied into $\alpha_t$ when semantic region labels are available; in the current experiments the control map is derived directly from the refined candidate-support map. In other words, $p_t$ remains the spatial support and $\alpha_t$ only modulates the perturbation amplitude within that support.

Combining Eq.~\eqref{eq:ncp} with Eq.~\eqref{eq:image_noise} gives the
effective spatial weight
\begin{equation}
\label{eq:effective_weight}
w_t = \alpha_t \odot p_t
    = \alpha\,\frac{p_t^2}{\max(p_t)+\epsilon}.
\end{equation}
The ``penalty'' is therefore a normalization and amplitude-control operation,
not an additive training loss; $w_t$ is the weight the protocol redistributes.

We use a Gaussian-inspired calibration index to select the base image-space noise level:
\begin{equation}
\label{eq:gaussian_calibration}
\sigma_0(\varepsilon,\delta) =
\frac{C\sqrt{2\log(1.25/\delta)}}{\varepsilon},
\end{equation}
where $C$ is a user-selected image-space scale, not a validated sensitivity
bound. Equation~\eqref{eq:gaussian_calibration} is a calibration index for
perturbation strength and is \emph{not} sufficient to claim end-to-end
differential privacy: no adjacency relation, release sensitivity bound, or
composition proof is supplied. Here $\varepsilon$ and $\delta$ are index
parameters only; we fix $\delta=10^{-5}$ and choose $C$ so that $\sigma_0$ is
in pixel units, making $\sigma_0=8$ the fixed operating point of the retrieval
runs rather than a certified guarantee. Noise is drawn deterministically per
frame from a fresh seed offset, not accumulated across frames.

\subsection{Selective Image-Space Perturbation}
After DCRF and NCP, PPEDCRF perturbs only the inferred candidate background support:
\begin{equation}
\label{eq:image_noise}
I_t' = \mathrm{clip}\!\left(I_t + \alpha_t \odot p_t \odot \eta_t,\ 0,\ 255\right),
\qquad
\eta_t \sim \mathcal{N}(0,\sigma_0^2 I).
\end{equation}
Equation~\eqref{eq:image_noise} is the released mechanism evaluated in this paper. The implementation indexes each Gaussian draw by the release seed and frame index. The attacker never observes intermediate maps such as $u_t$, $p_t$, or $\alpha_t$; it only observes the sanitized frame $I_t'$ and computes its own embedding $f(I_t')$ for retrieval. The product $\alpha_t \odot p_t$ is intentional: $p_t$ defines where perturbation is allowed, while $\alpha_t$ further increases the amplitude on the highest-confidence sensitive pixels inside that same support.

Selective perturbation matters because global Gaussian noise can damage foreground perception disproportionately, while random masking perturbs many pixels that are irrelevant to location retrieval. PPEDCRF instead concentrates perturbation energy on a candidate background support; temporal reuse is a consistency prior, not a sequence-level privacy guarantee.

\begin{comment}
\section{Experimental Evaluation}
\label{sec:experiments}

\subsection{Experimental Protocol}
\textbf{Controlled paired-scene retrieval benchmark.}
To evaluate the background-based retrieval attacker reproducibly, we build a controlled benchmark from a frozen pool of 120 monitoring sequences, taking one representative frame each. A fixed ResNet18~\cite{he2016deep} embedder greedily pairs the most similar sequences into 12 synthetic location pairs, with the gallery view drawn from the first half of one sequence and the query view from the second half of its partner; 36 unused hard distractors, selected by maximum anchor similarity, give galleries of size 12, 24, and 48. Location identifiers are synthetic: the pool has no GPS labels, so this benchmark does not measure recovery of a verified geographic identity. Construction is deterministic and all stochastic results average three noise seeds.

The benchmark is intentionally hard: anchor cosine similarity averages
$\approx$0.99 across the 12 pairs and distractors average 0.86 maximum
similarity to the paired set, so we treat it as a hard-negative proxy.

Downstream utility is checked on the same sanitized frames against 200 indexed
VOC segmentation masks and a co-audited 200-image COCO detection manifest,
with frozen DeepLabV3--ResNet50 and Faster R-CNN ResNet50-FPN; both manifests
were audited together so the reported mIoU and mAP share one checkpoint, code
path and build. The values are in the Supplementary Material and agree with
the retrieval-side ablation: at this budget the mechanism preserves more
downstream utility than global noise or the deterministic baselines, the
module ablations leave it unchanged at this rounding, and mask-guided mosaic
loses the most. Legacy detector-training results on MOT16/17, Cityscapes,
KITTI and VOC2008 exist in the released repository but reflect a training-time
pipeline unrelated to the location-retrieval question and are not used.

\begin{comment}
\begin{table}[t]
\centering
\caption{E4 same-image utility on 200 indexed VOC segmentation images and 200 COCO detection images at $\sigma_0=8$. The original row is the unsanitized reference; all other rows report metrics on the exact sanitized frames. mAP@50 uses a frozen Faster R-CNN ResNet50-FPN; mIoU uses a frozen DeepLabV3-ResNet50.}
\label{tab:e4seg}
\begin{tabular}{lcc}
\hline
Method & mAP@50 $\uparrow$ & mIoU $\uparrow$ \\
\hline
Original reference & 0.475 & 0.697 \\
PPEDCRF & 0.429 & 0.669 \\
w/o temporal consistency & 0.429 & 0.669 \\
w/o NCP & 0.429 & 0.669 \\
Unary-only + NCP & 0.429 & 0.669 \\
No DCRF & 0.428 & 0.669 \\
Mask-guided blur & 0.402 & 0.626 \\
Mask-guided mosaic & 0.195 & 0.296 \\
Global Gaussian noise & 0.352 & 0.632 \\
\hline
\end{tabular}
\end{table}
\end{comment}

Both agree with the retrieval-side ablation: at this budget the mechanism
preserves more downstream utility than global noise or the deterministic
baselines, the module ablations leave same-image utility unchanged at this
rounding --- mirroring their weak separation on the retrieval side --- and
mask-guided mosaic loses the most on both tasks, consistent with it being the
most visually destructive baseline in Table~\ref{tab:ablation}.

\begin{comment}
\begin{table}[t]
\centering
\caption{Datasets and roles in the revised evaluation.}
\label{tab:datasets}
\begin{tabular}{p{0.24\linewidth} p{0.68\linewidth}}
\hline
Dataset / benchmark & Role in this paper \\
\hline
MOT16 / MOT17, Cityscapes, KITTI, VOC2008 & Legacy utility references; not used in this paper. \\
VOC segmentation (E4) & Same-image downstream utility check on the indexed sanitized query frames. \\
Controlled paired-scene benchmark & New retrieval, ablation, fixed-budget privacy--utility, and attacker-sensitivity evaluation built from monitoring sequences. \\
\hline
\end{tabular}
\end{table}
\end{comment}

Unless otherwise stated, frames are resized to $192\times 320$, the attacker uses cosine similarity on fixed embeddings, and privacy is reported as Top-1/5/10 retrieval accuracy. The query is one released frame; neighboring frames feed only the optional temporal support estimate. For query $q_i$ we define the correct-versus-hardest-negative margin $m_i=s(q_i,g_i^+)-\max_{g_j\in\mathcal{G}_i^-}s(q_i,g_j^-)$, positive when the correct item outranks every negative. Perceptual utility uses PSNR and SSIM~\cite{wang2004image}. Raw queries are deterministic; sanitized variants average three noise seeds (1234--1236). All backbone-specific results come from one unified benchmark construction and the same averaging protocol, so attacker families are directly comparable.

\textbf{Scope of this benchmark.} The paired-scene benchmark isolates release-side privacy mechanisms when a full geo-localization dataset is unavailable locally. Anchoring pair mining on ResNet18 similarity improves reproducibility but may bias the benchmark toward retrieval models resembling the anchor embedder; the ResNet50/VGG16 and dedicated VPR transfers (CosPlace, MixVPR, Patch-NetVLAD) partly test this without eliminating it. Results therefore support claims about privacy--utility trade-offs, masking strategy, and attacker sensitivity, but do not prove universal robustness or constitute a formal differential-privacy certificate. Appendix~A provides implementation-alignment notes and a 50-pair confirmation assessing whether the 12-pair conclusions persist on a broader set.

\textbf{External validation on real place identities (E1).} To check that the
findings are not an artifact of mined proxy pairs, we built gate-passing
manifests from the official Mapillary Street-Level Sequences
(MSLS)~\cite{mapillarymsls2020} \texttt{train\_val} release, using its
official \texttt{unique\_cluster} place label. The primary manifest and the
two official cross-time subtasks were evaluated against six attacker
backbones --- ResNet50, VGG16, and the VPR architectures
CosPlace~\cite{berton2022rethinking}, MixVPR~\cite{ali2023mixvpr}, and
Patch-NetVLAD~\cite{hausler2021patchnetvlad}, alongside ResNet18 --- first at
200 queries over a 1000-image two-city gallery, then rebuilt at 8-city,
400-query, 2,000-gallery scale. On real place data the protective effect is
directionally consistent with the proxy but far smaller (Top-1 drop 0.018
versus 0.111); it holds in all six backbones on the wider primary manifest,
while 2 of 12 cross-time cells remain adverse, both ResNet18. A place-cluster
bootstrap finds only 5 of the original 18 cells individually significant, so
we treat cell counts as directional. Queries remain almost entirely
Forward-view and day, with empty season/weather fields. A white-box
sign-gradient diagnostic --- a distinct adaptive-attacker threat model, not
pooled with the black-box results --- collapses Top-1 to $0.000$--$0.055$ on
every backbone and manifest scale tested. Appendix~A reports all cell-level
comparisons, coverage gates, and scope.

\textbf{Independent check on the support map itself.} Separately from the
retrieval protocol, we asked whether the learned map agrees with an
independent notion of image sensitivity. Retraining a mask-backed predictor on
KITTI-360 (semantic labels as a structural-support target, not sensitivity
ground truth) and testing on a held-out sequence yields non-constant maps but
negative attribution agreement: Spearman $-0.00024$, top-10 overlap 0.0497,
and 0.0947 of the occlusion-reference energy captured. This is an independent
negative result consistent with the energy-matched finding below; Appendix~A
gives the protocol.

\textbf{Comparability across reported tables.} The main retrieval table uses the proxy12 export with three noise seeds, the default ResNet18 attacker, and fixed $\sigma_0=8$; the full cross-backbone/gallery matrix, the proxy50 confirmation, and the matched-PSNR sweep are in Supplementary Material. Because each proxy run shares pair discovery, gallery construction, and seeds across variants, within-run comparisons are directly meaningful.

\subsection{First Axis: Does Placement Matter?}
\label{sec:causality}

The protocol's first output on this mechanism is a diagnostic about the
mechanism itself, and it is worth reporting before any retrieval number. The
released support map is, numerically, a constant. Across inputs its unary
sigmoid ranges over $[0.4991, 0.5015]$ with a spatial coefficient of variation
of $6\times10^{-4}$, and the effective weight $w_t$
(Eq.~\eqref{eq:effective_weight}) inherits that flatness (support coverage
$1.000$: every pixel is ``selected''). The mechanism marketed as spatially
selective performs no spatial selection at all; at this checkpoint it is
global Gaussian noise scaled by one half.

This matters for how the redistribution controls must be read. Spreading,
rolling, or permuting an already-flat map returns an essentially identical
map, so those three controls are near-no-ops here, and the zero-discordant
outcome they produce follows from the map's degeneracy rather than from any
fact about placement. We report this explicitly because the opposite reading
--- treating that null as evidence that location is irrelevant --- would be
unsupported, and because it is precisely the kind of error the protocol
exists to prevent. No privacy--utility curve would have revealed this: the
mechanism still produces a perfectly plausible trade-off plot while its map
does nothing.

Testing placement therefore requires placements that are genuinely
non-uniform and independent of the degenerate map. We construct six: the
learned map, a uniform control, spectral-residual saliency, Sobel edge
magnitude, a centre-bias prior, and a fixed random field, each renormalised
to carry exactly the learned map's sum of squared weights (energy-conservation
gate, worst relative error $3.4\times10^{-5}$). We add two more from the
attacker's own input gradients: an \emph{oracle} that concentrates the budget
on the pixels where $|\partial\,\mathrm{cos}(f(x),g^{+})/\partial x|$ is
largest, and its anti-oracle. If placement can matter at all, the oracle is
where it should show.

A degenerate map, however, also makes a weak testbed: one could object that
the learned placement was never given a chance. We therefore repeat the entire
study on a second checkpoint of the same architecture, trained with structural
mask supervision, whose map is genuinely non-uniform (coefficient of variation
$0.158$, versus $6\times10^{-4}$ for the released one). Both studies use 6
backbones, 3 seeds, and both the 12-pair and 50-pair benchmarks: 2,928 paired
rows and 49 paired comparisons against the uniform control each.

Table~\ref{tab:placement} reports both. The result is consistent and
one-directional: \emph{across all 98 comparisons, no placement is ever
significantly better than uniform}. Several are significantly worse. With the
degenerate map, edge placement raises Top-1 by $0.067$ on average,
fixed-random by $0.051$, and the attacker-gradient oracle by $0.028$; none of
these reaches cluster-robust significance. With the genuinely selective map
the picture sharpens, and against the mechanism: the learned placement itself
raises Top-1 by $0.039$ on average and is \emph{significantly} worse than
uniform on VGG16 ($+0.194$, 95\% CI $[0.028, 0.361]$), as is saliency
placement on the 50-pair run ($+0.053$, CI $[0.007, 0.113]$). Giving the map
real spatial selectivity did not help privacy; where it had a measurable
effect, it hurt.

The attacker-gradient oracle is the most informative row. Even when it is
allowed to concentrate the entire budget on the pixels the attacker is most
sensitive to, it lands within noise of uniform ($+0.028$ and $+0.008$ on the
two checkpoints), and its anti-oracle --- which deliberately avoids those
pixels --- performs the same ($-0.001$, $-0.003$). At this budget, targeting
the right pixels and pointedly avoiding them are indistinguishable.

\subsection{The Null Holds on Real Geographic Data}
\label{sec:msls_placement}

Everything above is measured on the mined paired-scene corpus, whose
``locations'' are synthetic pair identifiers. That is the obvious place for
the result to be an artifact, so we repeat the whole study against the real
place-labelled MSLS benchmark: 400 queries drawn from eight cities, a shared
2,000-image gallery, 277 official \texttt{unique\_cluster} place identities,
and correctness decided by place label rather than by pair identity. Each
comparison now rests on 1,200 paired observations, roughly an order of
magnitude more than the proxy affords, and the energy-conservation gate passes
at $5.6\times10^{-7}$.

The null survives intact. Against the uniform reference, no placement differs
significantly in any of the seven comparisons, and every difference is inside
$\pm0.01$ Top-1: the attacker-gradient oracle reaches $-0.0008$ ($p=1.00$),
the learned map, fixed-random and saliency all $+0.0008$, the anti-oracle
$+0.0025$, edge $+0.0050$, and centre bias $+0.0092$ ($p=0.19$). The
mechanism's own map is again indistinguishable from spreading the budget
uniformly, and the oracle again fails to convert its knowledge of the
attacker into retrieval privacy.

\textbf{Placements driven by published models.} Every placement above is
either the mechanism's own or a prior we wrote, which is a fair objection to
the generality of the result. We therefore add placements that neither we nor
the mechanism designed, taken from off-the-shelf segmentation models trained
by others. Each implements the strategy this class of mechanism claims to
follow --- perturb the background scene, spare the foreground objects --- but
reads the scene through a different published model, so agreement between them
is evidence about the strategy rather than about one network. Two of them,
DeepLabV3-ResNet50~\cite{chen2017rethinking} and
FCN-ResNet50~\cite{long2015fully}, identify the scene as the COCO/VOC
\texttt{background} class, the complement of the object categories. The third,
SegFormer-B0~\cite{xie2021segformer} trained on ADE20K~\cite{zhou2017scene},
instead names the built environment directly,
and we sum the probabilities of its thirteen scene-structure classes
(building, sky, road, tree, sidewalk, wall, house, skyscraper, grass, plant,
earth, path, fence). The two taxonomies are genuinely different readings of
the same frame: over 60 real query frames their maps correlate $0.69$ on
average, ranging from $-0.33$ to $0.98$, so this is not one rule
re-parameterised three times. All three
discriminate spatially where the released checkpoint does not, with
coefficient of variation $0.07$--$0.49$ against its $6\times10^{-4}$.

None of them beats uniform. Against the uniform reference within the same run,
over 1,200 paired observations each: DeepLabV3 $\Delta=-0.0008$ (95\% CI
$[-0.010, +0.008]$), FCN-ResNet50 $\Delta=+0.0008$ ($[-0.007, +0.008]$), and
SegFormer/ADE20K $\Delta=0.0000$ ($[-0.009, +0.009]$), every one with exact
McNemar $p=1.00$; the mechanism's own learned map sits at $+0.0008$ alongside
them. Per-query energy agrees with the uniform control to a maximum relative
error of $7\times10^{-7}$ and mean PSNR to $0.02$\,dB, so the comparison is
between placements and not between budgets.

\textbf{Why we do not simply re-run a published system.} The stronger test
would put an existing published mechanism through the protocol, and we
searched for one. The requirement is specific: the method must emit a
per-pixel placement decision, so its weights can be renormalised to a common
energy, and it must have released code. We found none meeting both for this
threat model. Recent work on geolocation privacy optimises a full-image
adversarial perturbation rather than deciding where to place a budget --- the
direction-optimising setting our white-box baseline already covers, not a
placement rule --- while a 2024 survey of image privacy protection organises
the field by protection level, treats neither background nor geolocation
privacy as a category, and identifies no spatially selective method with
released code; generative anonymisers replace content instead of allocating a
budget, so they have no weight map to redistribute.

We regard this as a finding rather than an excuse: the control we propose is
absent from this literature partly because the literature has not published
systems in a form where it could be applied. What we can do instead, and do,
is instantiate the placement rule with off-the-shelf published models rather
than priors of our own design.

The sensitivity statistics replicate almost exactly, which is what makes the
mechanistic account portable rather than dataset-specific. On real imagery the
attacker's per-pixel sensitivity has coefficient of variation $1.01$ (against
$1.12$ on the proxy) with the top decile carrying $64.7\%$ of gradient energy
(against $67.7\%$), while the learned map's own top decile captures $9.4\%$
(against $9.9\%$) --- chance --- at rank correlation $0.020$ (against
$0.009$). Sensitive pixels are just as concentrated on real geographic data,
the learned map is just as blind to them, and targeting them is just as
ineffective.

\subsection{Placement Matters, But Only Outside the Operating
Regime}
\label{sec:budget}

Everything above is measured at the mechanism's stated operating point,
$\sigma_0=8$. Since a null result invites the objection that the budget was
simply too small for placement to express itself, we repeat the study across
$\sigma_0 \in \{4, 16, 32, 50\}$ on both checkpoints. The budget is doing real
work over this range: uniform Top-1 falls monotonically from $0.778$ at
$\sigma_0=4$ to $0.444$ and $0.361$ at $\sigma_0=50$ on the two checkpoints.

Up to and including the operating point, the null holds: over 56 further
comparisons at $\sigma_0 \in \{4, 16\}$ and above, no placement is
significantly better than uniform. At the two largest budgets, however, the
picture changes, and we report it because it qualifies our own thesis. Running
$\sigma_0 \in \{32, 50\}$ at the 50-pair scale (150 paired observations per
comparison) makes edge-magnitude placement \emph{significantly better} than
uniform in three independent cells: $\Delta=-0.140$ on the released
checkpoint at $\sigma_0=50$ (33 discordant pairs, $p=3\times10^{-4}$, 95\% CI
$[-0.247,-0.040]$), and $-0.087$ and $-0.133$ on the mask-supervised
checkpoint at $\sigma_0=32$ and $50$ (CIs $[-0.173,-0.007]$ and
$[-0.253,-0.013]$). The attacker-gradient oracle turns consistently negative
over the same range ($-0.040$ to $-0.060$) without reaching significance.

So placement is not inert in general. It becomes exploitable once the
perturbation is large enough, and what it must align with is image structure
--- edges --- rather than the ``sensitive region'' a learned map nominates.
The practical consequence is unfavourable for this class of design: the regime
where placement starts to pay is the regime where the frame has already been
degraded far past the quality these mechanisms exist to preserve. At
$\sigma_0=50$ the attacker's Top-1 has already halved on its own. Selectivity
claims are made in the high-quality regime, and that is precisely the regime
where our data show placement cannot help. Across every budget tested, the
mechanism's own learned map never beats uniform ($\pm0.000$, $+0.047$,
$+0.007$).

\begin{table}[t]
\centering
\caption{Energy-matched placement rules on two checkpoints of the same mechanism: the released one, whose map is numerically constant, and a mask-supervised one whose map is genuinely selective. $\Delta$ is the mean paired Top-1 difference against the uniform control over 6 backbones, 3 seeds and both benchmarks; \emph{positive $\Delta$ means worse privacy}. $\ast$ marks the only cluster-robust significant results, both of which are harmful. No placement is significantly better than uniform in any of the 98 comparisons.}
\label{tab:placement}
\begin{tabular}{lcc}
\hline
 & \multicolumn{2}{c}{$\Delta$ vs uniform} \\
Placement & constant map & selective map \\
\hline
anti-oracle (gradient) & $-0.001$ & $-0.003$ \\
learned support        & $\pm0.000$ & $+0.039^{\ast}$ \\
oracle (gradient)      & $+0.028$ & $+0.008$ \\
saliency               & $+0.026$ & $+0.031^{\ast}$ \\
centre bias            & $+0.030$ & $+0.030$ \\
fixed random           & $+0.051$ & $+0.039$ \\
edge magnitude         & $+0.067$ & $+0.052$ \\
\hline
\end{tabular}
\end{table}

\subsection{The Operator Is the Live Variable}
\label{sec:operator}

Everything above varies where a budget is spent while holding fixed what it is
spent on: additive isotropic Gaussian noise, which is the operator this
literature almost universally uses. That leaves the null open to a reading we
cannot dismiss from the placement data alone --- perhaps placement is inert
only because this particular operator is. We therefore add the missing axis
and vary the operator instead, holding the placement fixed.

Doing so requires a stronger matching rule than the one used so far.
Comparing operators at equal \emph{nominal} weight energy is meaningless when
their distortion profiles differ by construction, so every condition in this
section is released at the MSE the uniform Gaussian reference delivers on the
same frames, solved per frame by bisection through the pixel clamp. This also
retires a confound we disclose in Section~\ref{sec:mechanism}: matching
delivered distortion rather than nominal weight energy removes the clipping
asymmetry outright instead of merely bounding it.

\textbf{In the controlled model.} With the Jacobian known and delivered energy
identical to within $10^{-5}$, three operators that differ in \emph{spatial}
structure are interchangeable. Against a clean Top-1 of $0.867$, uniform
placement gives $0.832$ under isotropic noise, $0.843$ under spatially
correlated noise, and $0.845$ under a deterministic-magnitude random-sign
perturbation; what placement buys on top is likewise unchanged ($-0.43$,
$-0.42$, $-0.48$). A fourth operator, identical in energy but with its sign
chosen along the discriminative direction, gives $0.001$ --- with uniform
placement, that is, with no spatial selectivity whatsoever, and placement adds
nothing further because there is nothing left to remove. Spatial structure in
the operator is worth as little as spatial structure in the map; directional
control is worth everything.

\textbf{On real imagery.} We repeat this with four operators on the real MSLS
benchmark, each placed by the most concentrated rule available (edge,
top-decile share $0.839$) and by the uniform control, at matched delivered
MSE (Table~\ref{tab:operator}). At the operating point the null survives every
one of them: no edge-versus-uniform difference is significant and none exceeds
$0.01$ in magnitude. The negative result is a property of the regime, not of
the operator we happened to start with.

At a budget eight times larger, two things appear that the placement-only view
cannot see. First, the operator now matters a great deal: at identical
delivered distortion and identical uniform placement, Top-1 ranges from
$0.159$ under additive noise down to $0.053$ under block quantisation, a
threefold spread bought purely by changing what the budget is spent on.
Second, the placement effect that finally emerges is \emph{operator-dependent
in sign}: concentrating on edges helps under additive noise and hurts under
low-pass and block quantisation. A placement rule tuned on one operator can
therefore be actively harmful on another.

\begin{table}[t]
\centering
\caption{Perturbation operators at matched delivered MSE on the real
place-labelled benchmark, 1,200 paired observations per arm. ``Uniform'' is
Top-1 under the uniform control; $\Delta$ is edge placement minus uniform,
with \emph{negative meaning better privacy}. $\ast$ marks exact-McNemar
significance at $\alpha=0.05$. At the operating point no operator makes
placement matter; at the larger budget the operator changes Top-1 threefold
and the sign of $\Delta$ reverses between operators.}
\label{tab:operator}
\begin{tabular}{llccc}
\hline
MSE & Operator & Uniform $\downarrow$ & $\Delta$ (edge) & $p$ \\
\hline
15.7 & additive Gaussian    & 0.1950 & $+0.0025$ & 0.82 \\
15.7 & correlated Gaussian  & 0.1892 & $-0.0008$ & 1.00 \\
15.7 & selective low-pass   & 0.1975 & $+0.0100$ & 0.18 \\
15.7 & block quantisation   & 0.2000 & $-0.0075$ & 0.31 \\
\hline
241.5 & additive Gaussian   & 0.1592 & $-0.0433^{\ast}$ & $<$0.001 \\
241.5 & correlated Gaussian & 0.0650 & $+0.0167$ & 0.067 \\
241.5 & selective low-pass  & 0.1025 & $+0.0500^{\ast}$ & $<$0.001 \\
241.5 & block quantisation  & 0.0525 & $+0.1150^{\ast}$ & $<$0.001 \\
\hline
\end{tabular}
\end{table}

\textbf{What this says for design.} Read together, the two axes give a
recommendation with a sign rather than only a warning. At the quality-preserving
budgets these systems target, no placement rule beats uniform under any
operator we tested, whereas the choice of operator moves retrieval accuracy by
a factor of three at matched distortion once the budget is large enough to move
anything at all. The literature's usual procedure --- fix the operator at
additive Gaussian noise and optimise the map --- optimises the variable that
does not matter while holding fixed the one that does. Our own testbed is an
instance of that procedure, which is why its map turning out to be constant
cost it so little.

This account also predicts the budget dependence in
Section~\ref{sec:budget}: Eq.~\eqref{eq:displacement} is a small-perturbation
statement, so as amplitude grows the response becomes structure-dependent and
placement acquires leverage --- which is why edge-aligned placement, not
sensitivity-map placement, is what starts to win at $\sigma_0 \ge 32$. Retrieval privacy in
the high-quality regime is therefore governed by alignment with the
embedding's sensitive directions, and spatial selectivity is the wrong control
variable exactly where these mechanisms operate.

\subsection{Why Allocation Is the Wrong Variable}
\label{sec:mechanism}

Two results now need one explanation: placement does nothing at the operating
point under any operator, and the operator itself only starts to matter once
the budget is large. Both follow from the same observation, that a fixed
budget of zero-mean perturbation controls how far an embedding moves but not
which way, and only the second of those decides a ranking.

\textbf{What first-order theory predicts.} It is worth being precise about
what should happen, because the prediction is the opposite of what we
measure. Write the attacker embedding as $f$ and the released frame as
$x+\delta$, with $\delta_i = w_i\varepsilon_i$, $\varepsilon_i$ i.i.d.\
$\mathcal{N}(0,\sigma^2)$, and the placement fixed by the weights $w$ subject
to $\sum_i w_i^2 = E$. For small $\delta$, $\Delta f \approx J\delta$ with
$J=\partial f/\partial x$, so
\begin{equation}
\label{eq:displacement}
\mathbb{E}\lVert \Delta f\rVert^2 \;\approx\; \sigma^2 \sum_i w_i^2\,\lVert J_i\rVert^2 ,
\end{equation}
where $J_i$ is the column of $J$ for pixel $i$. Under the fixed-energy
constraint, the right-hand side is maximised by putting all the weight on the
pixels of largest $\lVert J_i\rVert$ --- exactly the oracle placement. Since
we measured that $\lVert J_i\rVert$ is strongly concentrated (top decile
carrying 67.7\% of the energy), Eq.~\eqref{eq:displacement} predicts a
substantial oracle advantage. There is none.

\textbf{Why the prediction fails.} Two effects break it, and both are
measurable. First, the released frame is bounded. The protocol matches
$\operatorname{mean}(w^2)$ exactly, which absent clipping would force the
delivered mean squared error to match exactly as well; any observed shortfall
is therefore energy lost at the clamp. There is such a shortfall, and it grows
with concentration: on the proxy, saliency delivers MSE 15.13 and edge 15.33
against uniform's 15.70, up to 3.6\%, while on the real benchmark, where the
placements we could compare are far less concentrated, it falls to 0.25\%.
Notably the clipping \emph{fraction} is identical across placements
($1.91\%$ on the real benchmark), so the loss is in how much energy each
boundary pixel sheds, not in how many pixels reach the boundary --- a
distinction worth stating because the fraction alone would suggest, wrongly,
that clipping is not involved. Second, and more fundamentally,
Eq.~\eqref{eq:displacement} governs the \emph{magnitude} of the embedding
displacement, whereas retrieval is decided by rank order. A larger
$\lVert\Delta f\rVert$ in a direction that does not reorder the gallery
changes no outcome, and for isotropic $\varepsilon$ the direction of
$J\delta$ is essentially arbitrary within the column space of $J$ regardless
of where $\delta$ is placed.

The second point is the one that generalises, and it is what the white-box
contrast isolates: an optimized perturbation at the same PSNR chooses the
direction rather than the location, and it succeeds completely.

\textbf{A caveat on our own protocol, and its removal.} Matching energy in
\emph{weight} space to a relative error of $10^{-7}$ does not match delivered
distortion: concentrated placements receive an effectively smaller budget. The
confound runs in the direction of making concentrated placements look worse,
so it never rescued them --- they already lose, and on the real benchmark the
shortfall is only $0.25\%$ --- but it is a real defect of nominal-energy
matching. Section~\ref{sec:operator} matches achieved MSE instead, solved per
frame through the clamp, and the conclusion is unchanged under either rule. We
state the confound rather than leave it for a reader to find.

\textbf{How much selectivity is actually being expressed.} A null across
placement rules invites the objection that the rules were never far from
uniform to begin with, so we measured how concentrated each one is. Writing
the top-decile energy share as the fraction of total squared weight carried by
the 10\% of pixels with the largest weights --- exactly $0.100$ for uniform by
construction --- the rules span most of the available range: edge $0.839$
(coefficient of variation $1.77$), centre $0.649$, saliency $0.353$, fixed
random $0.242$, and the three published segmentation models $0.113$, $0.115$
and $0.126$.

Two things follow, and they cut in opposite directions. The segmentation
placements are only mildly selective, and the reason is worth stating plainly:
on street imagery almost every pixel \emph{is} scene structure, so ``perturb
the background, spare the foreground'' is close to a no-op there --- the
strategy has little to work with in the setting where location privacy
actually matters. That is a limitation of the strategy, not a limitation of
our test of it. But the objection does not extend to the study as a whole,
because edge placement concentrates $84\%$ of the budget into a tenth of the
pixels, an eightfold departure from uniform, and still does not beat uniform
at the operating point. Strong selectivity was available and tested; it simply
did not pay until the budget grew large enough to destroy the image anyway.

The oracle's failure is the informative part, because it rules out the
obvious explanation. One might expect placement to be irrelevant because the
attacker's sensitivity is spatially homogeneous, leaving nothing to target.
We measured this and it is false. Per-pixel sensitivity is strongly
concentrated: its coefficient of variation is $1.12$ (range $0.84$--$2.56$),
and the top decile of pixels carries $67.7\%$ of the total gradient energy.
There is a well-defined set of pixels that matter to the attacker.

The learned map is simply blind to it. Its own top decile captures $9.9\%$ of
the gradient energy --- indistinguishable from the $10\%$ a random decile
would capture --- and its rank correlation with attacker sensitivity is
$0.009$ (range $[-0.028, 0.034]$). This is consistent with the independent
attribution diagnostic reported above, and unsurprising given that the map is
constant.

Yet the oracle, which by construction places the entire budget on exactly
those high-sensitivity pixels, still does not beat uniform --- for the two
reasons derived above.

\textbf{The strictly correct oracle.} That oracle targets the similarity to
the correct gallery item, but Top-1 is decided by the \emph{margin} between
that similarity and the best competitor: a perturbation that lowers both
equally reorders nothing. The oracle that maximises the right objective
therefore places the budget by the margin gradient, and we ran it. It is
indeed the better signal --- its point estimate, $-0.0083$, is an order of
magnitude larger than the similarity oracle's $-0.0008$ and in the favourable
direction --- but over 1,200 paired observations it does not clear the noise
floor ($p=0.275$, 95\% CI $[-0.027, +0.009]$). Placing the budget by exactly
the quantity that decides the outcome, with white-box access to compute it,
still buys nothing measurable. The objection that we simply built the wrong
oracle is thereby closed.

\subsection{A Controlled Check with an Exactly Known Jacobian}
\label{sec:known-jacobian}

On real data the account above rests on an estimated Jacobian, so it remains
an interpretation. We therefore repeat the question in a setting where nothing
has to be estimated: a retrieval task whose encoder is known in closed form,
so per-pixel sensitivity is exact and every quantity in
Eq.~\eqref{eq:displacement} is computable. The encoder's column norms are
given a heavy-tailed profile matched to the concentration measured on the real
attacker (top decile carrying $67.7\%$ of sensitivity energy). Queries and
their positives are two different draws of the same latent place, as in real
place recognition, and a nuisance parameter sets how far apart those two views
are and therefore how hard the clean task is. Each cell averages three seeds
and 600 trials.

The first result is a sanity check that the algebra is right, and it is:
measured mean squared displacement matches Eq.~\eqref{eq:displacement} to
within $4\%$ across all 168 cells. The second is that the predicted oracle
advantage is not merely theoretical. When the clean task is easy, the oracle
is devastating --- at clean Top-1 $0.86$ it drives retrieval down by $0.42$
relative to the uniform control at identical energy. So the null we measure on
images is \emph{not} a consequence of the first-order argument being wrong.

Three candidate explanations for the null can then be switched on
individually. Bounding the released signal roughly halves the advantage at
every operating point ($-0.42 \rightarrow -0.33$, $-0.19 \rightarrow -0.07$),
confirming that concentrated placements shed energy at the clamp rather than
delivering it. Making the clean retrieval task harder decays it monotonically:
$-0.42$ at clean Top-1 $0.86$, $-0.15$ at $0.34$, $-0.04$ at $0.12$, and
$+0.00$ at $0.03$. Replacing the linear encoder with a nonlinear one, so its
Jacobian is valid only locally, changes nothing ($-0.37$ at clean $0.70$,
$-0.02$ at clean $0.07$) --- which rules out linearisation error and leaves
clipping and task difficulty as the operative causes. A budget sweep in the
same model reproduces the reversal structure we measure on images: the
advantage grows with $\sigma$, from $+0.02$ at the smallest budget to $-0.22$
at large budgets.

\textbf{Where the controlled model and the measurement disagree.} It would be
convenient to stop there, but the calibration does not fully close. Our real
benchmark has clean Top-1 $0.210$, which the operating point reduces to
$0.195$ --- a drop of $0.015$. Matching that drop inside the controlled model
puts it at a setting where the model still predicts an oracle advantage of
about $-0.15$, whereas the measurement is $-0.0008$ with a 95\% interval of
$\pm 0.01$. The idealised oracle therefore outperforms the real one by more
than an order of magnitude.

We report this rather than tune the model until it agrees, because the
direction of the discrepancy is itself informative. The controlled oracle is
handed the exact Jacobian; the real one gets a single gradient evaluated at
the clean frame, which is a weaker guide to where a finite perturbation will
actually move the embedding. Our real oracle is in fact consistently
\emph{worse} than uniform across the budget sweep ($+0.028$ at $\sigma_0=16$,
$+0.111$ at $32$, $+0.084$ at $50$), so on images an attacker-informed
placement does not merely fail to help --- it preserves retrieval slightly
better than spreading the budget. The controlled study thus bounds how much a
placement rule could win with perfect information, and the gap between that
bound and what we measure is the price of estimating sensitivity from a
gradient. Either way the practical conclusion is unchanged, and if anything
strengthened: placement is not where the privacy comes from at the budgets
these mechanisms operate at.

The contrast with the white-box attacker reported above closes the argument.
A sign-gradient perturbation at a \emph{comparable PSNR} drives Top-1 to
$0.000$--$0.055$, because it optimizes the perturbation's direction rather
than its location. Location without direction achieves nothing; direction without any spatial
restriction achieves everything.

\subsection{Controlled Paired-Scene Retrieval Benchmark}
Table~\ref{tab:ablation} reports the proxy12 default setting: ResNet18 attacker, gallery size 48, and $\sigma_0=8$. The raw query baseline reaches Top-1 0.833. The mechanism changes this to $0.722\pm0.048$ while preserving 36.17\,dB PSNR and 0.945 SSIM, whereas global Gaussian noise reaches a lower Top-1 ($0.639\pm0.127$) at only 30.20\,dB PSNR and 0.827 SSIM; the deterministic baselines go further still, to 0.583 at 28.08\,dB (blur) and 0.417 at 26.07\,dB (mosaic). The ordering is monotone in distortion: every variant that reduces retrieval
more also degrades the image more. Top-10 saturates at 1.000 here, so we
report Top-1 and Top-5.

The module ablations say the same thing: removing temporal consistency,
replacing NCP with a fixed control map or dropping DCRF leaves rounded
retrieval and quality unchanged. With Section~\ref{sec:causality} this is the
same finding from another angle --- at fixed energy the mechanism's internal
structure has no measurable effect on what the attacker retrieves.

\begin{table}[t]
\centering
\caption{Proxy12 paired-scene benchmark, default setting: ResNet18 attacker, gallery size 48, $\sigma_0=8$, 12 paired locations, and 36 hard distractors. Lower retrieval accuracy indicates stronger privacy. Stochastic variants report mean$\pm$std over three noise seeds; deterministic baselines are exact.}
\label{tab:ablation}
\resizebox{\columnwidth}{!}{%
\begin{tabular}{lcccc}
\hline
Method & Top-1 $\downarrow$ & Top-5 $\downarrow$ & PSNR (dB) $\uparrow$ & SSIM $\uparrow$ \\
\hline
Raw query & 0.833 & 1.000 & -- & -- \\
PPEDCRF & $0.722\pm0.048$ & 1.000 & 36.17 & 0.945 \\
w/o temporal consistency & $0.722\pm0.048$ & 1.000 & 36.17 & 0.944 \\
w/o NCP (fixed control) & $0.722\pm0.048$ & 1.000 & 36.17 & 0.944 \\
Unary-only + NCP & $0.722\pm0.048$ & 1.000 & 36.17 & 0.945 \\
No DCRF & $0.722\pm0.048$ & 1.000 & 36.17 & 0.944 \\
Mask-guided blur & 0.583 & 0.917 & 28.08 & 0.897 \\
Mask-guided mosaic & 0.417 & 0.833 & 26.07 & 0.858 \\
Global Gaussian noise & $0.639\pm0.127$ & $0.972\pm0.048$ & 30.20 & 0.827 \\
\hline
\end{tabular}}
\end{table}

\subsection{Matched Quality: the Mechanism Against One Other Operator}
\label{sec:matched}

The fixed-budget comparison of Table~\ref{tab:ablation} is a comparison at
unequal distortion: the mechanism is spending less of it and buying less
privacy with it. Equalising quality removes the difference entirely. Over a 12-point matched-PSNR sweep the mechanism, its four ablations
and global noise have identical displayed Top-1 at all three targets, with no
discordant outcome in 15 comparisons ($n=36$, exact McNemar $p=1.0$,
query-cluster bootstrap 95\% CI $[0.000,0.000]$); extending to five further
backbones and to gallery sizes 12 and 100 leaves one non-significant
discordant pair in 120 comparisons. This is a non-rejection rather than proven
equivalence --- six discordant pairs would be needed at $\alpha=0.05$.

In the vocabulary of Section~\ref{sec:operator} this is the mechanism measured
against a single alternative operator, global Gaussian noise, at matched
quality; the deterministic baselines are compared the same way there, across
four operators and at matched delivered MSE rather than matched PSNR. The
$5.97$\,dB quality advantage the mechanism shows at equal nominal $\sigma_0$
is not a selectivity benefit either: with a map constant at $\approx0.5$ it is
simply applying half the amplitude, and the quality follows from that scaling.
What the mechanism delivers is a favourable \emph{quality} position at a given
nominal noise scale --- a real engineering property, but not the privacy
property it was designed for.

\begin{comment}
\begin{figure*}[!t]
    \centering
    \includegraphics[width=0.8\textwidth]{figs/privacy_utility_tradeoff.jpg}
    \caption{Fixed-budget proxy12 comparison at $\sigma_0=8$ for the ResNet18 attacker and gallery size 48. Lower retrieval accuracy means stronger privacy; higher PSNR means higher visual utility. Error bars show the three-seed retrieval standard deviation when applicable.}
    \Description{Two scatter plots comparing Top-1 and Top-5 retrieval accuracy with PSNR for PPEDCRF, global Gaussian noise, mask-guided blur, and mask-guided mosaic at a fixed noise budget.}
    \label{fig:tradeoff}
\end{figure*}
\end{comment}

The proxy12 transfer audit varies gallery size and attacker backbone over
three seeds and eight backbones; the per-cell values are in the released
artifact. Under ResNet18 the mechanism moves Top-1 from 0.833 to 0.750 at
gallery 12 and to $0.722\pm0.048$ at galleries 24 and 48. Across the 24
backbone--gallery cells, 20 are negative, one is tied and three are positive,
with all three positive cells belonging to ResNet50; VGG16, both CLIP
variants, CosPlace and Patch-NetVLAD are negative in every condition, and
CosPlace is stable at $-0.167$ throughout. The proxy50 scaling check is more
favourable but still not universal: 23 of 24 cells negative, one tied, none
positive. Transfer is therefore bounded and attacker-dependent rather than
uniform --- the ResNet50 cells are why it should not be presented as
attacker-agnostic --- and because both runs are mined paired-scene proxies
rather than GPS-grounded benchmarks, they are evidence of that bounded effect
and not a universal guarantee.

\begin{comment}


\begin{figure*}[!t]
    \centering
    \begin{subfigure}[t]{0.95\textwidth}
        \centering
        \includegraphics[width=\linewidth]{figs/retrieval_robustness_topk_top.jpg}
        \caption{Top row: ResNet18, ResNet50, VGG16, and CLIP ViT-B/32.}
    \end{subfigure}

    \vspace{0.4em}

    \begin{subfigure}[t]{0.95\textwidth}
        \centering
        \includegraphics[width=\linewidth]{figs/retrieval_robustness_topk_bottom.jpg}
        \caption{Bottom row: CLIP ViT-L/14, CosPlace, MixVPR, and Patch-NetVLAD.}
    \end{subfigure}

    \caption{Attacker-backbone and gallery-size sensitivity from the proxy12 eight-backbone benchmark (three seeds), rendered as two stacked four-panel subfigures. PPEDCRF reduces Top-1 in 20 of 24 cells, ties one cell, and increases it in three ResNet50 cells.}
    \Description{Two stacked four-panel line charts. The top subfigure shows ResNet18, ResNet50, VGG16, and CLIP ViT-B/32. The bottom subfigure shows CLIP ViT-L/14, CosPlace, MixVPR, and Patch-NetVLAD. Each panel compares raw and PPEDCRF Top-1 retrieval accuracy across gallery sizes 12, 24, and 48, with seed-averaged confidence intervals.}
    \label{fig:robustness}
\end{figure*}
\end{comment}

The robustness analysis exposes important patterns for interpretation. First, CosPlace, MixVPR, and Patch-NetVLAD are supportive in most proxy12 gallery conditions, while ResNet50 is adverse in every condition. Second, both CLIP variants have negative deltas in all three galleries, but their absolute changes are smaller than those of CosPlace or Patch-NetVLAD. Third, the benchmark remains rank-sensitive rather than fully obfuscating: Top-5 and Top-10 often remain high, and stronger attacker-adaptive perturbations would still be needed for more decisive privacy guarantees.

\subsection{Sequence-Length Attacker and Temporal Consistency}
\label{sec:sequence}
A stronger multi-frame attacker does not change the frame-level picture. Over
clip lengths 1--7 and four pooling strategies (50 pairs, 100-candidate
gallery, 6 backbones, 3 seeds; 33,600 rows; Appendix~A), Top-1 stays within
$[0.669,0.683]$, close to raw ($[0.670,0.682]$), and stronger pooling does not
close the already-small gap. Global noise remains more disruptive
($[0.608,0.616]$) at lower quality. The temporal term does deliver what it was
designed for --- flicker falls from $7.11\pm5.78$ to $3.87\pm4.78$ --- which,
consistent with the rest of this paper, is a stability benefit rather than a
privacy one; we claim no sequence-level guarantee.

\begin{comment}
\begin{figure*}[!t]
    \centering
    \includegraphics[width=\textwidth]{figs/qualitative_figure.jpg}
    \caption{Qualitative visualization of PPEDCRF sanitization ($\sigma_0=24$ for visual clarity; at lower $\sigma_0$ the difference map shows sharper spatial selectivity). (a)~Original video frame. (b)~DCRF candidate-support heatmap overlay, where warm colors indicate larger predicted support values in the background. (c)~Noise-protected output. (d)~Absolute pixel difference showing spatially selective perturbation. (e)~Mask-guided blur on the same support. (f)~Zoomed crop comparing original (left) and protected (right) in a selected background-support region.}
    \Description{Six-panel qualitative visualization: original frame, DCRF candidate-support heatmap overlay, noise-protected frame, pixel difference map, mask-guided blur result, and zoomed side-by-side crop of a selected background-support region.}
    \label{fig:sidebyside}
\end{figure*}
\end{comment}

\section{Conclusion}
A selective visual privacy mechanism makes two choices: where to spend a
perturbation budget, and what to spend it on. The literature learns, tunes and
ablates the first, and inherits the second without comment. We built a
protocol that controls both --- redistributing a mechanism's own weights under
a verified energy-conservation gate, and comparing operators at matched
delivered distortion --- and the two axes behave nothing alike.

Placement is inert. Across eight placement rules, three rules driven by
published segmentation models, an attacker-gradient oracle, the strictly
correct margin-gradient oracle, two checkpoints, six attacker backbones, four
operators and a real place-labelled benchmark, nothing beats spreading the
budget uniformly at the quality-preserving operating point; the map with
genuine spatial selectivity is significantly \emph{worse} on one backbone. The
protocol also exposed that our testbed's released map is numerically constant,
so it performs no spatial selection at all while still producing a plausible
privacy--utility curve --- a defect invisible to the evaluation such systems
normally receive, and one that cost the mechanism almost nothing precisely
because placement was never carrying the protection.

The operator is not inert, but it is not a free lunch either. At the operating
point every operator we tested leaves retrieval roughly where uniform Gaussian
noise leaves it. Once the budget is large enough to move retrieval at all,
changing only what it is spent on moves Top-1 from $0.159$ to $0.053$ at
identical delivered distortion and identical uniform placement, and the
placement effect that finally appears reverses sign between operators --- so a
placement rule tuned on one operator is actively harmful on another.

A controlled retrieval task with an exactly known Jacobian says why. At
matched energy, three operators differing in \emph{spatial} structure are
interchangeable, while one differing in \emph{directional} control drives
Top-1 to $0.001$ under uniform placement and no selectivity whatsoever. What
buys privacy is alignment with the embedding's sensitive directions, which is
also why an optimized white-box perturbation succeeds at comparable quality
where every placement rule fails. Spatial allocation --- of the map or of the
operator --- is the wrong control variable, and it is wrong in the regime
these systems are built for.

We do not claim a well-chosen operator solves the problem: none of ours does
at quality-preserving budgets, and directional control in the strong sense
requires attacker knowledge a deployed defense will not have. The claim is
comparative and procedural. Among the choices a designer actually makes, the
operator carries leverage and the placement does not, so a claim that a
learned map improves privacy should be reported alongside an energy-matched
placement control and a matched-distortion operator comparison. Both are
cheap, and either can overturn the claim.


\appendices
\section{Released Implementation Notes}
The injector perturbs the \emph{released image}, not a hidden feature tensor.
The unary predictor is a newly implemented encoder--decoder, not a pretrained
VPR or segmentation model: two $3\times3$ convolution blocks at 32 and 64
channels with max-pooling, two transposed-convolution upsampling blocks and a
$1\times1$ convolution to one logit channel, 87,441 parameters in total,
trained with AdamW at $10^{-4}$ and binary cross-entropy where masks are
available. Layer-by-layer definitions are in the released artifact.

Neither checkpoint validates a privacy-sensitivity map. The released one
records a null mask-root fallback and emits near-constant probabilities; the
mask-backed KITTI-360 diagnostic below closes the execution gate but still
finds near-zero attribution agreement.

For exact DCRF reproduction, the sigmoid in Eq.~(2) is the element-wise
$\sigma(z)=1/(1+\exp(-z))$, the spatial smoother is average pooling with a
$3\times3$ kernel, stride 1, and reflect padding, and the reported benchmark
uses five mean-field-style iterations with spatial and temporal weights both
equal to 2. Optical-flow warping is disabled; the previous refined map is
reused directly when temporal refinement is enabled.

The normalized control penalty (NCP) is also implemented conservatively in the released code. Rather than building a large symbolic hierarchy at runtime, the current implementation normalizes the refined support map and uses it as a per-pixel perturbation controller. The released sanitizer intentionally multiplies the refined support map and the normalized control map, so high-confidence selected pixels receive the strongest perturbation inside the same spatial support. The manuscript therefore treats the hierarchy-based interpretation as motivation, while the empirical sections report results only for the released normalization rule.
More explicitly, the effective weight in the released sanitizer is
\begin{equation}
\label{eq:appendix_effective_weight}
w_t=\alpha_t\odot p_t
    =\alpha\,\frac{p_t^2}{\max(p_t)+\epsilon},
\end{equation}
which makes clear that NCP is a normalization and amplitude-control rule,
not an additional additive penalty term.


\section{Energy-Matched Redistribution Controls}
\label{sec:spatial_causality}

This section documents the three redistribution controls and the diagnostic
they produced. Let $w_t=\alpha_t\odot p_t$ be the effective spatial weight. We
compare \texttt{full} (the unmodified DCRF/NCP output) against
\texttt{uniform\_energy} (every pixel at the amplitude giving the same mean
squared weight), \texttt{rolled\_energy} (the map cyclically shifted by a
seed- and frame-dependent offset) and \texttt{permuted\_energy} (the same
values reassigned to random positions), evaluated at $\sigma_0=8$, gallery
sizes 50/75/100, six backbones, 50 paired locations and three seeds
(11,700 rows).

\textbf{Energy-conservation gate.} Each control must preserve the mean squared
weight of \texttt{full} to a relative tolerance of $10^{-6}$. Realised errors
are far below it (worst case $4.76\times10^{-7}$), and the gate passes for
every row.

\textbf{Result and its limit.} All three controls reproduce \texttt{full}'s
Top-1 outcome exactly: zero discordant pairs and exact McNemar $p=1.0$ for
each. This null cannot on its own establish that placement is irrelevant,
because the released checkpoint's map is numerically constant --- unary output
spanning $[0.4991, 0.5015]$, spatial coefficient of variation
$6\times10^{-4}$, and \texttt{support\_coverage} $1.0$ on every row, meaning
every pixel is ``selected''. Redistributing an already-flat map returns
essentially the same map, so the controls are near-no-ops here. We report the
null as the protocol's first diagnostic --- the mechanism performs no spatial
selection while still producing a plausible privacy--utility curve --- and
establish the placement claim itself in Section~\ref{sec:causality}, on
placements that do not depend on this map.

\textbf{Perturbation footprint.} Over the 2,700 \texttt{full} rows the mean
squared effective weight is $0.2506$ (std $4.1\times10^{-6}$), the mean
effective MSE is $15.70$ (std $0.175$), and $1.54\%$ of output pixels hit the
$[0,255]$ clamp (std $0.92\%$), identical across backbones since clipping is a
property of the image rather than the attacker.

\section{Placement-Rule Study}
\label{sec:placement_appendix}

This section reports the full data behind the placement comparison. Because
the released checkpoint's map is constant, the redistribution controls cannot
settle whether location matters; these placements are constructed
independently of the learned map.

\textbf{Placements.} Each rule produces a non-negative per-pixel map, rescaled
so its sum of squared weights equals the learned map's for the same clip,
seed and frame, so outcome differences are attributable to location rather
than magnitude. Besides the learned map and a uniform reference we evaluate
spectral-residual saliency, Sobel gradient magnitude, a centre-bias Gaussian
prior and a fixed smooth random field, plus two placements derived from the
attacker itself: an \emph{oracle} weighting each pixel by
$|\partial\,\mathrm{cos}(f(x),g^{+})/\partial x|$ and an \emph{anti-oracle}
inverting it.

\textbf{Protocol.} Two checkpoints are evaluated: the released one, whose map
is constant (spatial coefficient of variation $6\times10^{-4}$), and a
mask-supervised one of the same architecture whose map is genuinely
non-uniform (coefficient of variation $0.158$). Each is run over 6 attacker
backbones, 3 noise seeds, and both the 12-pair (gallery 48) and 50-pair
(gallery 100) benchmarks, giving 2,928 paired rows and 49 paired comparisons
against the uniform reference per checkpoint. The energy-conservation gate
passes with worst-case relative error $3.4\times10^{-5}$. Pairing is within a
run: query identifiers are synthetic and are regenerated per benchmark, so
rows are keyed by their source run to prevent the two benchmarks from being
paired against each other.

\begin{table}[t]
\centering
\caption{Placement-rule study. $\Delta$ is the mean paired Top-1 difference against the uniform reference; positive means \emph{worse} privacy. ``Zero-disc.'' counts comparisons (of 7 run--backbone cells) with no discordant pairs. No placement is significantly better than uniform in either checkpoint.}
\label{tab:placement_full}
\resizebox{\columnwidth}{!}{%
\begin{tabular}{lcccccc}
\hline
 & \multicolumn{3}{c}{constant map} & \multicolumn{3}{c}{selective map} \\
Placement & Top-1 & $\Delta$ & Zero-disc. & Top-1 & $\Delta$ & Zero-disc. \\
\hline
uniform (ref.)   & 0.636 & --- & --- & 0.634 & --- & --- \\
anti-oracle      & 0.635 & $-0.001$ & 6/7 & 0.631 & $-0.003$ & 5/7 \\
learned          & 0.636 & $\pm0.000$ & 7/7 & 0.673 & $+0.039$ & 2/7 \\
oracle           & 0.664 & $+0.028$ & 0/7 & 0.642 & $+0.008$ & 0/7 \\
saliency         & 0.661 & $+0.026$ & 2/7 & 0.665 & $+0.031$ & 2/7 \\
centre bias      & 0.666 & $+0.030$ & 1/7 & 0.664 & $+0.030$ & 1/7 \\
fixed random     & 0.686 & $+0.051$ & 1/7 & 0.673 & $+0.039$ & 1/7 \\
edge magnitude   & 0.703 & $+0.067$ & 0/7 & 0.686 & $+0.052$ & 0/7 \\
\hline
\end{tabular}%
}
\end{table}

\textbf{Significance.} Of the 98 comparisons, none favours a placement over
uniform at the cluster-robust level. On the constant map, one comparison
reaches exact-McNemar significance (edge/CosPlace, $+0.167$, 6 discordant
pairs, $p=0.031$) but its query-cluster interval includes zero. On the
selective map, two comparisons are cluster-robust significant and both are
harmful: the learned placement on VGG16 ($+0.194$, 95\% CI
$[0.028, 0.361]$) and saliency on the 50-pair run ($+0.053$, CI
$[0.007, 0.113]$). We report the exact test and the cluster bootstrap
separately throughout, and let the cluster-robust verdict govern where they
disagree, because the three seeds of a query are not independent draws.

\textbf{Attacker sensitivity.} The statistics behind
Section~\ref{sec:mechanism} are computed here. Beyond the figures quoted
there, the mask-supervised map is more variable than the released one but no
better aligned on average (top-decile capture $10.3\%$, rank correlation
$0.013$), though individual frames range from $1.1\%$ to $47.1\%$.

\subsection{Budget Dependence}
\label{sec:placement_budget}

The placement results above are measured at $\sigma_0=8$. To establish where
that null holds and where it breaks, we sweep
$\sigma_0 \in \{4, 16, 32, 50\}$ on both checkpoints (ResNet18, 12 paired
locations, gallery 48, 3 seeds), and repeat the two largest budgets at the
50-pair scale with gallery 100 for adequate power. Energy-conservation gates
pass throughout (worst relative error $4.0\times10^{-7}$).

The budget is doing real work across this range. Uniform Top-1 falls
monotonically from $0.778$ at $\sigma_0=4$ to $0.444$ (released checkpoint)
and $0.361$ (mask-supervised) at $\sigma_0=50$, so the sweep spans a regime
where the perturbation goes from mild to severe.

\begin{table}[t]
\centering
\caption{Placement $\Delta$ versus uniform across budgets (ResNet18, 12 pairs, gallery 48). Positive is worse privacy. Edge placement reverses sign as the budget grows; the mechanism's own learned map never improves on uniform at any budget.}
\label{tab:placement_budget}
\resizebox{\columnwidth}{!}{%
\begin{tabular}{lcccccccc}
\hline
 & \multicolumn{4}{c}{constant map} & \multicolumn{4}{c}{selective map} \\
Placement & $\sigma_0{=}4$ & 16 & 32 & 50 & $\sigma_0{=}4$ & 16 & 32 & 50 \\
\hline
learned      & $0.000$ & $0.000$ & $0.000$ & $0.000$ & $-0.028$ & $0.000$ & $+0.111$ & $+0.028$ \\
anti-oracle  & $0.000$ & $-0.028$ & $0.000$ & $0.000$ & $0.000$ & $0.000$ & $+0.028$ & $0.000$ \\
oracle       & $-0.028$ & $+0.028$ & $+0.111$ & $+0.083$ & $-0.028$ & $+0.111$ & $+0.139$ & $+0.056$ \\
saliency     & $0.000$ & $+0.056$ & $+0.056$ & $+0.056$ & $0.000$ & $+0.083$ & $+0.111$ & $+0.056$ \\
centre bias  & $+0.028$ & $+0.139$ & $+0.167$ & $+0.139$ & $+0.028$ & $+0.167$ & $+0.222$ & $+0.222$ \\
fixed random & $+0.028$ & $+0.083$ & $+0.111$ & $+0.056$ & $+0.028$ & $+0.111$ & $+0.111$ & $+0.028$ \\
edge         & $+0.028$ & $+0.083$ & $-0.111$ & $-0.167$ & $+0.028$ & $+0.083$ & $-0.111$ & $-0.139$ \\
\hline
\end{tabular}%
}
\end{table}

\textbf{High-budget confirmation at scale.} The sign reversal for edge
placement is confirmed at 50 pairs (150 paired observations per comparison),
where it becomes the only placement to beat uniform at the cluster-robust
level, in three independent cells: $\Delta=-0.140$ on the released checkpoint
at $\sigma_0=50$ (33 discordant pairs, exact $p=3.2\times10^{-4}$, 95\% CI
$[-0.247,-0.040]$); $-0.087$ on the mask-supervised checkpoint at
$\sigma_0=32$ (23 discordant, $p=0.011$, CI $[-0.173,-0.007]$); and $-0.133$
at $\sigma_0=50$ (36 discordant, $p=0.0012$, CI $[-0.253,-0.013]$). The
attacker-gradient oracle moves consistently in the same direction over this
range ($-0.040$ to $-0.060$) without reaching significance, while centre-bias
and saliency remain harmful.

We report this because it bounds our own claim. Placement is exploitable, but
only after the perturbation has grown well beyond the quality these mechanisms
are designed to preserve --- at $\sigma_0=50$ the attacker's Top-1 has already
fallen by roughly a third on the uniform baseline alone --- and what pays
there is alignment with image edges rather than with a learned sensitivity
map, which never improves on uniform at any budget tested.


\section{Matched-PSNR/Effective-MSE Comparison}
\label{sec:matched_psnr}

The fixed-budget comparison reports each variant at its own operating point,
confounding visual quality with retrieval privacy. To separate them we run the
proxy12 protocol (ResNet18 attacker, gallery 48, three seeds) at twelve noise
scales $\sigma\in\{4,\dots,50\}$ for every sigma-tunable variant and select,
per target PSNR, the swept point of closest achieved PSNR. Since
$\operatorname{PSNR}=10\log_{10}(255^2/\operatorname{MSE})$ for 8-bit frames
at fixed dimensions, matching PSNR matches output-distortion energy up to
rounding, though not the latent draws or supports. Table~\ref{tab:matched_psnr}
gives three targets; the six sigma-tunable variants are collapsed into one row
each because their Top-1 is identical to three decimals at every target, which
is the result. The deterministic baselines sit at lower PSNR than any swept
target yet reach lower Top-1, so over the reachable quality range the two
families trade privacy against utility differently; they are compared properly,
at matched delivered MSE, in Section~\ref{sec:operator}.

\begin{table}[t]
\centering
\caption{Matched-PSNR/effective-MSE comparison, ResNet18 attacker, gallery size 48, proxy12 protocol, three seeds. The five sigma-tunable variants (PPEDCRF and its four ablations) and global Gaussian noise are collapsed into one row per target because their Top-1 is identical to three decimals at every target; the matched $\sigma$ and achieved PSNR ranges across those six are given. Mask-guided blur/mosaic have no sigma parameter here and are reported at their single native point with the resulting PSNR gap. Top-1 $\downarrow$: lower means better retrieval privacy.}
\label{tab:matched_psnr}
\begin{tabular}{lccc}
\hline
Target / variant & Matched $\sigma$ & Actual PSNR & Top-1 $\downarrow$ \\
\hline
\multicolumn{4}{l}{\emph{30\,dB target}} \\
All six sigma-tunable & 8--16 & 30.18--30.20\,dB & 0.639 \\
Mask-guided blur & --- & 28.08\,dB & 0.583 \\
Mask-guided mosaic & --- & 26.07\,dB & 0.417 \\
\hline
\multicolumn{4}{l}{\emph{33\,dB target}} \\
All six sigma-tunable & 6--12 & 32.66--32.68\,dB & 0.667 \\
Mask-guided blur & --- & 28.08\,dB & 0.583 \\
Mask-guided mosaic & --- & 26.07\,dB & 0.417 \\
\hline
\multicolumn{4}{l}{\emph{36\,dB target}} \\
All six sigma-tunable & 4--8 & 36.17--36.19\,dB & 0.722 \\
Mask-guided blur & --- & 28.08\,dB & 0.583 \\
Mask-guided mosaic & --- & 26.07\,dB & 0.417 \\
\hline
\end{tabular}
\end{table}

The prespecified paired analysis for each of the five non-reference variants
against global noise at each target (15 comparisons) finds $n=36$ paired
query--seed outcomes, zero discordant pairs, exact McNemar $p=1.0$ and a
query-cluster bootstrap 95\% CI of $[0.000,0.000]$ throughout. This is a
non-rejection rather than proven equivalence: at $n=36$ the smallest
discordant asymmetry reaching $\alpha=0.05$ is six pairs.

\subsection{Deterministic Baseline Matched-PSNR Comparison}
\label{sec:deterministic_matched_psnr}

An earlier version of this comparison swept the blur kernel and mosaic block
size to bring the deterministic baselines to the sigma-tunable variants' PSNR
targets. Because the kernel is a discrete odd integer, no setting landed
closer than $1.60$\,dB to the 36\,dB target, so the comparison could only be
reported with an explicit residual gap. Section~\ref{sec:operator} supersedes
it: treating blur and block quantisation as \emph{operators} and solving per
frame for the gain that hits the reference MSE matches distortion exactly
rather than approximately, on the real place-labelled benchmark, with paired
significance over 1,200 observations per arm. The conclusion there is the
stronger form of what this sweep indicated --- the deterministic operators buy
privacy by spending image quality, and at matched distortion they are more
protective than additive noise only once the budget is already large.


\bibliographystyle{IEEEtran}
\bibliography{ref}

\end{document}
